\subsubsection{Codificación de pares tríticos y suma con acarreo}
\label{mg04:codificacion}

La lectura ternaria de las emisiones APP--TRIT--TPK tiene un ambiente
de seis coordenadas sobre \(\mathbb F_3\). Agruparlas en tres pares
ordena ese ambiente mediante tres dígitos nonádicos. La agrupación
conserva la posición y el orden de los trits; la suma que representa
debe transportar también su acarreo.

Sea \(T=\{0,1,2\}\) el conjunto de representantes ternarios.
Para una clase \(x\in\mathbb Z/9\mathbb Z\), escribimos
\(\widetilde x\in\{0,\ldots,8\}\) para su representante canónico.
Definimos
\begin{equation}
 \beta:T^2\longrightarrow\mathbb Z/9\mathbb Z,
 \qquad \beta(a,b)=[a+3b]_9.
 \label{mg04:beta}
\end{equation}
Aquí \([s]_m\) designa, cuando se emplea como dígito, el representante
en \(\{0,\ldots,m-1\}\) de la clase de \(s\).

\begin{proposition}[Inversa y transporte de la suma]
El mapa \(\beta\) es una biyección, con inversa
\begin{equation}
 \beta^{-1}(x)=
 \left([\widetilde x]_3,
       \left\lfloor\frac{\widetilde x}{3}\right\rfloor\right).
 \label{mg04:inversa}
\end{equation}
La ley sobre \(T^2\) transportada desde la suma nonádica es
\begin{equation}
 (a,b)\oplus(c,d)=
 \left([a+c]_3,
 \left[b+d+\left\lfloor\frac{a+c}{3}\right\rfloor\right]_3\right).
 \label{mg04:suma}
\end{equation}
Con esta ley, \(\beta\) es un isomorfismo de grupos aditivos.
El neutro es \((0,0)\) y la inversa de \((a,b)\) es
\begin{equation}
 \left([-a]_3,
 \left[-b+\left\lfloor\frac{-a}{3}\right\rfloor\right]_3\right).
 \label{mg04:opuesto}
\end{equation}
\end{proposition}

\begin{proof}
La división euclídea de cada \(0\leq\widetilde x\leq8\) por
tres da un único resto \(a\in T\) y un único cociente \(b\in T\),
con \(\widetilde x=a+3b\). Esto prueba la biyección y
\eqref{mg04:inversa}.

Para la suma, escribamos \(a+c=3k+r\), con
\(r=[a+c]_3\) y \(k=\lfloor(a+c)/3\rfloor\). Entonces
\[
 (a+3b)+(c+3d)=r+3(b+d+k).
\]
La reducción de \(b+d+k\) módulo tres sólo cambia el miembro
derecho en un múltiplo de nueve. Por tanto,
\(\beta((a,b)\oplus(c,d))=\beta(a,b)+\beta(c,d)\).
La biyectividad transporta asociatividad, conmutatividad y neutro.
Por último, la división \(-a=3\lfloor-a/3\rfloor+[-a]_3\)
aplicada a \(-a-3b\) da \eqref{mg04:opuesto}.
\end{proof}

El acarreo es visible ya en
\[
 (2,0)\oplus(1,0)=(0,1),\qquad
 \beta(2,0)+\beta(1,0)=[3]_9.
\]
La suma vectorial de \(\mathbb F_3^2\) produciría \((0,0)\)
en ese mismo ejemplo. Asimismo, \((1,0)\) tiene orden nueve para
\(\oplus\), mientras que su orden para la suma vectorial es tres.
Estas operaciones especifican las dos estructuras algebraicas que
conviven sobre el mismo conjunto de nueve representantes.

\begin{corollary}[Coordenatización cúbica del ambiente ternario y fibras de proyección]
La aplicación coordenada a coordenada
\begin{equation}
 \begin{split}
 \beta^{\times3}:T^6&\longrightarrow(\mathbb Z/9\mathbb Z)^3,\\
 (t_1,t_2,t_3,t_4,t_5,t_6)&\longmapsto
 \bigl(\beta(t_1,t_2),\beta(t_3,t_4),\beta(t_5,t_6)\bigr)
 \end{split}
 \label{mg04:cubo}
\end{equation}
es una biyección de conjuntos de \(729=3^6=9^3\) elementos. La
proyección \(\operatorname{pr}_{12}(X,Y,Z)=(X,Y)\) sobre las
\(81\) posiciones nonádicas tiene exactamente nueve elementos
por fibra.
\end{corollary}

\begin{proof}
El producto de las tres inversas de \eqref{mg04:inversa} invierte
\eqref{mg04:cubo}. Fijadas \(X,Y\), la tercera coordenada
recorre las nueve clases de \(\mathbb Z/9\mathbb Z\); de ahí
\(\operatorname{pr}_{12}^{-1}(X,Y)=
\{(X,Y,Z):Z\in\mathbb Z/9\mathbb Z\}\) y el cardinal afirmado.
Conservar \(Z\) junto con \(X,Y\) permite aplicar la inversa
completa y recuperar los seis trits ordenados.
\end{proof}

El dominio de esta cuenta es el ambiente completo de \(729\)
palabras. El catálogo de \(468\) emisiones decimales y su imagen
admisible de \(243\) palabras ternarias conservan sus propios mapas
emisores y sus restricciones. La restricción de \(\beta^{\times3}\)
a esa imagen sigue siendo inyectiva, mientras que los cardinales de
sus fibras de proyección se determinan por la intersección con la
fibra ambiente. La afirmación uniforme de nueve elementos pertenece
al cubo completo.

La coordenatización bidimensional resultante recupera una representación etiquetada
del soporte APP ya construido. Las hojas de suma y producto son funciones adicionales sobre
ese soporte, con los representantes y cocientes declarados en el núcleo.
La tercera coordenada conserva el dato omitido por la proyección
facial; las rutas y la memoria temporal permanecen en el estado TPK
del que procede la lectura. De este modo, la codificación finita
explicita la información que transporta y la distingue de la historia
completa a la que pertenece.
